\documentclass[final,5p,times,twocolumn]{elsarticle}

\usepackage{amssymb}
\usepackage{amsmath}
\usepackage{booktabs}
\usepackage{graphicx}
\usepackage{array}
\usepackage{multirow}
\usepackage{hyperref}
\usepackage{microtype}
\usepackage{xcolor}

% Revision highlighting: text changed in the 2026-09-21 revision is set in blue.
% Set \revhighlightfalse for the clean submission copy.
\newif\ifrevhighlight
\revhighlighttrue
\newcommand{\rev}[1]{\ifrevhighlight{\color{blue}#1}\else#1\fi}

\input{generated/results_macros.tex}

\journal{Physiological Measurement}

\begin{document}

\begin{frontmatter}

\title{How far can a subject-independent sEMG array be reduced? A prespecified non-inferiority test of cross-fitted attribution for electrode selection}

\author[inst1]{Minh-Tri Ta}
\author[inst2,inst3]{Hoang-Khang Phan}
\author[inst4]{Le-Giang Tran\corref{cor1}}
\cortext[cor1]{Corresponding author. Email: tranlegiang@iuh.edu.vn}

\affiliation[inst1]{organization={College of Engineering and Computer Science, VinUniversity},
            city={Hanoi},
            country={Vietnam}}
\affiliation[inst2]{organization={Department of Biomedical Engineering, Faculty of Applied Science, Ho Chi Minh City University of Technology (HCMUT)},
            addressline={268 Ly Thuong Kiet Street, District 10},
            city={Ho Chi Minh City},
            country={Vietnam}}
\affiliation[inst3]{organization={Vietnam National University Ho Chi Minh City},
            addressline={Linh Trung Ward, Thu Duc City},
            city={Ho Chi Minh City},
            country={Vietnam}}
\affiliation[inst4]{organization={Faculty of Electronics Technology, Industrial University of Ho Chi Minh City},
            city={Ho Chi Minh City},
            country={Vietnam}}

\begin{abstract}
\textit{Objective.} To test, under \rev{prespecified decision rules}, whether explainability-guided electrode selection can reduce a forearm sEMG array without meaningful loss of subject-independent gesture recognition.
\textit{Approach.} We used 43 participants of GRABMyo (16 monopolar forearm electrodes, 16 active gestures) and leave-one-subject-out evaluation, \rev{in which z-scoring, model fitting and channel ranking used only the other participants; a label-free count-feature threshold, set once from pooled probe recordings, was identical in every leave-one-out recomputation}. Electrodes were ranked by TreeSHAP attributions cross-fitted over participant-disjoint inner folds. A stability-regularised variant (SR-CFA) penalised attribution that varied across those folds. A frozen analysis plan fixed the participant as the statistical unit, a 0.02 Macro-F1 non-inferiority margin, 10,000-resample cluster bootstrap intervals and Holm correction\rev{; the detailed test families for the secondary hypotheses were specified after the results existed}.
\textit{Main results.} With all electrodes, XGBoost reached a Macro-F1 of \FullXGB\ (95\% CI \FullXGBLCB--\FullXGBUCB). The smallest non-inferior subset was $k^* = \KStar$ electrodes ($\bar\Delta = \DeltaKTwelve$, 95\% CI $\DeltaKTwelveLCB$ to $\DeltaKTwelveUCB$)\rev{, unchanged after Holm correction over the 15 tested subset sizes}. One fewer electrode was not non-inferior, and $k^*$ was \KStarMarginPtOne\ and \KStarMarginPtThree\ under the 0.01 and 0.03 sensitivity margins. Attribution-ranked subsets outperformed \rev{the mean of} random subsets at 4, 6, 8 and 12 electrodes (largest gain $\HTwoKSixDiff$ at six). Deletion with retraining \rev{supported the faithfulness of} the ranking. The stability penalty changed almost no rankings and gave no measurable benefit. \rev{In exploratory, uncorrected comparisons, no statistically detectable difference was observed between the attribution ranking and a data-free, evenly spread layout at 4, 8 and 12 electrodes.} \rev{In an exploratory replication added after the main analysis (putEMG, 44 participants, 24 electrodes), $k^* = \PutKStarPtTwo$ of 24 at the same margin ($\PutKStarHolmPtTwo$ after Holm correction).}
\textit{Significance.} \rev{For subject-independent, offline recognition with this two-ring monopolar montage, about a quarter of the electrodes could be removed within a prespecified engineering tolerance, but not more.} Cross-fitted attribution identifies which positions matter, \rev{but in exploratory comparisons its advantage over an evenly spread layout was detectable only at six electrodes.}
\end{abstract}

\begin{keyword}
Surface electromyography \sep gesture recognition \sep electrode selection \sep SHAP \sep cross-fitting \sep non-inferiority \sep leave-one-subject-out
\end{keyword}

\end{frontmatter}

\section{Introduction}
\label{sec:introduction}

Surface electromyography (sEMG) pattern recognition is the dominant approach to multifunction myoelectric control~\cite{hudgins1993new,englehart2003robust,scheme2011electromyogram,farina2014extraction}. Recognition generally improves with more recording sites, but every electrode adds cost, wiring, contact area and a point of failure in a wearable device~\cite{vujaklija2016new}. This makes it attractive to identify a reduced set of electrodes that preserves recognition performance~\cite{phinyomark2012feature,liu2014emg,altimemy2013classification,stango2015spatial,qu2021reduce}.

Three methodological problems limit what published channel-reduction results can tell a device designer. First, most reductions are selected and evaluated within the same users. \rev{User-dependent evaluation is appropriate for personalised systems, but it answers a different question from the one a device intended for many users must answer: a layout chosen for one person need not transfer to the next, and subject-independent recognition is generally harder than user-dependent recognition~\cite{phinyomark2018emg,cote2019deep}. Separately, if data from a test user influences scaling, ranking or tuning, the reported subset is optimistically biased.} Second, feature-attribution methods such as Shapley additive explanations~\cite{lundberg2017unified,lundberg2020local} are increasingly used to rank channels~\cite{salmi2025shapaad}, but attributions can be unstable across data splits~\cite{visani2022statistical,kelodjou2023shaping,xiang2026roshap}, and a ranking computed on the evaluation data is itself a form of leakage. Third, the chosen number of channels is usually read off an accuracy curve by eye, with no pre-stated tolerance and no statement of what ``no meaningful loss'' means~\cite{schuirmann1987comparison,walker2011understanding}. A margin chosen after the curve is seen cannot support a non-inferiority claim~\cite{campbell2018noninferiority}.

This study tests explainability-guided channel reduction under conditions designed to remove these three sources of optimism. \rev{Each held-out participant is excluded from z-scoring, model fitting and channel ranking.} Channel rankings are computed by cross-fitting: TreeSHAP attributions are always evaluated on participants the attributing model did not see. \rev{The primary decision rules were fixed in a statistical analysis plan before any sparse-channel result existed: the non-inferiority margin ($\delta = 0.02$ Macro-F1), the participant as the statistical unit, the bootstrap, Holm correction within each confirmatory family, and the definition of the smallest acceptable channel count $k^*$. The detailed test families for the secondary hypotheses were specified afterwards, and we report where the analysis departed from the plan (Section~\ref{sec:deviations}).} We also test a stability-regularised variant of the ranking (SR-CFA), which penalises channels whose attribution varies across disjoint groups of training participants.

The study is primarily a test, not a method proposal. We asked four prespecified questions on the 43-participant GRABMyo dataset~\cite{pradhan2022multiday}:
(H1) whether some $k < 16$ is non-inferior to all 16 channels;
(H2) whether attribution-ranked subsets beat random subsets of the same size;
(H3) whether the stability penalty improves on attribution magnitude alone;
(H4) whether the attribution is faithful, i.e.\ whether removing highly ranked channels hurts more than removing random or low-ranked ones.
We report every answer, including the negative ones. Exploratory analyses covered cross-session transfer, a bipolar re-derivation of the montage, and \rev{methodological replications on NinaPro DB5~\cite{pizzolato2017comparison} and on putEMG~\cite{kaczmarek2019putemg}, a 44-participant, 24-electrode dataset added after all other results existed}. These are reported as exploratory.

\section{Related Work}
\label{sec:related_work}

\paragraph{Channel reduction for myoelectric control}
Channel and feature reduction for sEMG pattern recognition has been approached with feature-ranking and wrapper methods~\cite{phinyomark2012feature,liu2014emg}, with channel selection on multi-channel and high-density arrays~\cite{altimemy2013classification,stango2015spatial}, and with sparse-representation methods~\cite{qu2021reduce}. Large-scale reviews of the field emphasise that most evaluations are user-dependent. They also note that population-level (cross-user) models behave differently from per-user models~\cite{phinyomark2018emg}. Deep cross-user models and transfer learning are one response to this gap~\cite{cote2019deep}. Our study sits at the other end of the spectrum: a conventional feature-based classifier, a strictly subject-independent protocol, and a question about hardware rather than model capacity.

\paragraph{Attribution-guided selection and attribution stability}
TreeSHAP computes Shapley-value attributions for tree ensembles~\cite{lundberg2017unified,lundberg2020local}. Attribution has been used to reduce channels in other biosignal tasks, for example EEG auditory attention decoding~\cite{salmi2025shapaad}. Several works show that attributions can change across resampling and neighbourhood choices~\cite{visani2022statistical,kelodjou2023shaping}. A recent preprint treats Shapley attributions as distributions across resamples and proposes a robustness-adjusted importance score~\cite{xiang2026roshap}. That idea is close to the stability penalty tested here, and we do not claim the penalty itself as new. The general idea of preferring features that are selected consistently across subsamples goes back to stability selection~\cite{meinshausen2010stability}. What this study adds is the evaluation design. Attribution is cross-fitted over participants inside each outer training fold. The stability penalty is tested as a prespecified hypothesis against magnitude alone, and the selected subsets are judged by a prespecified non-inferiority rule.

\paragraph{Faithfulness of attributions}
A ranking that looks plausible need not reflect what the model uses. Deletion-based tests remove the features an explanation marks as important and measure the effect~\cite{samek2017evaluating}. Retraining after removal (ROAR) avoids confounding the explanation with distribution shift of the unretrained model~\cite{hooker2019benchmark}. We use the retraining variant with channels as the unit.

\paragraph{Electrode shift and multi-day use}
Electrode displacement and day-to-day variation degrade myoelectric pattern recognition~\cite{hargrove2006effect,kaufmann2010fluctuating}, and adaptation methods have been proposed to counter the resulting covariate shift~\cite{vidovic2016improving,muceli2014extracting}. GRABMyo was recorded over three days~\cite{pradhan2022multiday}, which allows an exploratory cross-session analysis. However, the evaluation here is subject-independent throughout. Its cross-session contrasts therefore combine a session effect with the cross-user gap, and they are not a pure test of donning and doffing.

\section{Methods}
\label{sec:methods}

\begin{figure*}[t]
\centering
\includegraphics[width=\textwidth]{figures/fig1_workflow.pdf}
\caption{\rev{Study workflow. For each held-out participant $p$, z-scoring, model fitting and channel ranking use the other participants $R_{-p}$; attributions are computed on participant-disjoint inner folds of $R_{-p}$; $k$-electrode subsets are retrained and tested on $p$. The same workflow was applied separately to each dataset.}}
\label{fig:workflow}
\end{figure*}

\subsection{Datasets}

\paragraph{GRABMyo (primary)}
GRABMyo v1.1.0~\cite{pradhan2022multiday} contains recordings from \GRABMyoParticipants\ able-bodied participants (\GRABMyoMale\ male, \GRABMyoFemale\ female; age \GRABMyoAgeMean\ $\pm$ \GRABMyoAgeSD\ years) on \GRABMyoSessions\ different days. Signals were recorded with an EMGUSB2+ amplifier (OT Bioelettronica) at 2048~Hz. In each session participants performed 16 hand and wrist gestures plus rest, with 7 trials of 5~s each. We analysed the 16 forearm electrodes. These form two circumferential rings of 8 electrodes (F1--F8 and F9--F16), 2~cm apart along the forearm, with electrode $j$ of one ring aligned with electrode $j+8$ of the other. Mean circumference at the ring position was \GRABMyoCircMean\ $\pm$ \GRABMyoCircSD~cm, giving about \GRABMyoArcSpacing~cm between neighbouring electrodes in a ring. The recordings are monopolar, so the selection unit in the primary analysis is a single electrode. All \GRABMyoFilesVerified\ distributed files matched the dataset's published SHA-256 checksums. All \GRABMyoRecords\ records were processed, with \GRABMyoRecordsFailed\ failures. The primary task is the 16 active gestures (rest excluded), for which chance Macro-F1 is about 0.063.

\paragraph{NinaPro DB5 (exploratory methodological replication)}
DB5~\cite{pizzolato2017comparison} provides \DBFiveSubjects\ able-bodied subjects (\DBFiveMale\ male, \DBFiveFemale\ female; age \DBFiveAgeMean\ $\pm$ \DBFiveAgeSD\ years, from the files' metadata) recorded with two Myo armbands (16 channels, 200~Hz). We used exercise~1 (12 finger movements, 6 repetitions each), segmented with the refined \texttt{restimulus}/\texttt{rerepetition} labels. The DB5 values are 8-bit integer codes, not physical units. The spectra show no powerline peak, and their power extends to the Nyquist frequency. DB5 was analysed separately and never pooled with GRABMyo. Its channel numbers bear no anatomical correspondence to GRABMyo's, so only the behaviour of the procedure is compared across datasets.

\rev{\paragraph{putEMG (exploratory methodological replication, added post hoc)}
putEMG~\cite{kaczmarek2019putemg} provides 44 able-bodied participants recorded on two days at least one week apart (7 to 82 days in these files; median 14). A 24-electrode monopolar array forms three circumferential bands of 8 electrodes (1--8, 9--16, 17--24) on the forearm, sampled at 5120~Hz with a 12-bit converter. Each session contains three trajectories with 20 repetitions of each of 7 active gestures, labelled from video-verified ground truth. This analysis was added at the end of the study, after all GRABMyo and DB5 results existed; its design was recorded before any putEMG model was fitted, and it is exploratory. All 264 gesture files were readable and complete, and every participant had both sessions. As for DB5, the putEMG channel numbers have no correspondence with GRABMyo's.}

\subsection{Preprocessing and features}
GRABMyo signals were band-pass filtered (20--450~Hz, 4th-order zero-phase Butterworth) and notch filtered at 50~Hz ($Q = 30$). A spectral audit chose 50~Hz over 60~Hz: mean band-power ratio was 10.35 at 50~Hz and 0.62 at 60~Hz. Windows were 250~ms with 50\% overlap and never crossed trial boundaries. From each channel we extracted ten features: mean absolute value (MAV), root mean square, variance, waveform length, zero crossings (ZC), slope sign changes (SSC), integrated EMG (IEMG), Willison amplitude (WAMP), mean frequency and median frequency. This gives a 160-dimensional vector per window, \GRABMyoWindows\ windows in total and \GRABMyoActiveWindows\ from active gestures. The ZC, SSC and WAMP amplitude threshold was set to $10^{-2}$~mV, about 1.15 times the median absolute first difference of the \rev{unfiltered signal in 54 probe records from 6 participants (rounded to one significant figure). This label-free scalar was set once rather than per fold. To check whether this let a held-out participant influence the features, we recomputed the median in every outer fold without that participant (post hoc). The same rule gave $10^{-2}$~mV in \LOOMonoIdentical\ of \LOOMonoFolds\ folds (leave-one-out medians \LOOMonoMin--\LOOMonoMax${}\times 10^{-3}$~mV; only \LOOMonoInProbe\ folds held out a probe participant). Fold-wise estimation would therefore have produced identical features.} An earlier threshold 87 times smaller had saturated WAMP and was discarded. For DB5 we used a 10--95~Hz band-pass, no notch, 250~ms windows (50 samples) and a threshold of 7 ADC units, derived in the same way \rev{from all ten subjects; the leave-one-out recomputation gave the same value in \LOODBFiveIdentical\ of \LOODBFiveFolds\ folds}. \rev{For putEMG, raw values were converted to mV and processed as for GRABMyo (20--450~Hz band-pass, 50~Hz notch, 250~ms windows of 1280 samples, 50\% overlap), keeping only windows that lay entirely inside one active-gesture segment (idle and rejected samples excluded). Because raw putEMG differences are quantised at one converter step, which equals their median, the threshold rule was applied to the filtered signal; it gave \PutThresh${}\times 10^{-3}$~mV and the same value in \PutThreshIdentical\ of 44 leave-one-out folds. This gives 240 features per window.}

\subsection{Classifiers and leave-one-subject-out evaluation}
All evaluation is leave-one-subject-out (LOSO): 43 folds for GRABMyo and 10 for DB5. In fold $p$, z-scoring statistics, model fitting and channel ranking use only $R_{-p}$. \rev{Apart from the label-free threshold checked above and the training-set sizes described below, the test participant contributes only to the score.} XGBoost used 300 trees, depth 6, learning rate 0.1, row and column subsampling of 0.8, and the \texttt{hist} method. It was trained on all \FoldTrainingWindows\ outer-training windows. For compute reasons, the RBF support vector machine ($C = 1$, $\gamma$ = \texttt{scale}) used a stratified 40,000-window subsample and the random forest (200 trees, minimum leaf size 2) a 200,000-window subsample. On one calibration fold, halving the SVM subsample from 40,000 to 20,000 windows lowered Macro-F1 from \SVMCalibFortyK\ to \SVMCalibTwentyK\rev{, suggesting that the SVM had not saturated; its reported result is therefore likely a lower bound. Model hyperparameters were fixed in advance and not tuned. The training-set sizes (the SVM and random-forest subsamples and the 200,000-window cap used for the channel analyses) were chosen from held-out Macro-F1 on two calibration folds (participants 1 and 7) before the confirmatory runs, a departure from strict fold-wise isolation for those two participants (Section~\ref{sec:deviations}). XGBoost used the \texttt{multi:softprob} objective and library defaults for all unlisted parameters.} The primary metric is participant-level Macro-F1 over the 16 active gestures; accuracy and balanced accuracy were also recorded. XGBoost was prespecified as the primary model for the channel analyses on grounds of compute.

\subsection{Cross-fitted channel attribution and the SR-CFA score}
Within each outer fold, $R_{-p}$ was split into $K = 5$ participant-disjoint inner folds. For inner fold $j$, an XGBoost model was trained on the other inner folds (capped at a \rev{uniformly sampled} 200,000-window subsample and z-scored with inner-training statistics only). Exact TreeSHAP attributions~\cite{lundberg2020local} were then computed for up to 20,000 windows of the inner-held-out participants\rev{, sampled uniformly without replacement from their pooled windows. The attribution is therefore window-weighted; because every participant contributes a similar number of windows (median \TestWindowsPerParticipant), this is close to participant-weighted}. Per-window channel importance averaged absolute attributions over the 16 classes and the 10 features of the channel,
\begin{equation}
\Phi_{i,c} = \frac{1}{G F}\sum_{g=1}^{G}\sum_{f=1}^{F} \lvert \phi_{i,g,c,f} \rvert ,
\end{equation}
and was averaged over windows and normalised to sum to one across channels, giving $I_{j,c}$. Across the five inner folds, $A_c = \operatorname{median}_j I_{j,c}$ measures attribution magnitude. $U_c = \operatorname{median}_j \lvert I_{j,c} - A_c \rvert$ measures its instability across disjoint groups of training participants. Channels were ranked \rev{(ties broken by lower channel index)} by
\begin{equation}
Q_c(\lambda) = A_c - \lambda\, U_c , \qquad \lambda \in \{0, 0.25, 0.5, 1\}.
\label{eq:srcfa}
\end{equation}
$\lambda = 0$ is cross-fitted SHAP magnitude; $\lambda > 0$ is the stability-regularised score (SR-CFA). A leakage assertion verified in every fold that no inner-held-out participant appeared in its inner training set. \rev{Unit tests verified that the outer test participant is never attributed, that inner folds are participant-disjoint, and that a synthetic dataset with a participant-confounded signal detects scaler leakage.}

The frozen plan required $\lambda$ to be chosen on outer-training data only but did not specify a criterion, and no selection procedure was implemented. We therefore select no $\lambda$ and report all four (a post-hoc specification, Section~\ref{sec:deviations}).

\subsection{Sparse sensor curve and non-inferiority (H1)}
For every outer fold, $\lambda$ and $k \in \{1,\ldots,16\}$, the top-$k$ channels of the fold's own ranking were used to retrain XGBoost from scratch. Each model used the fold's fixed 200,000-window stratified outer-training subsample and was tested on the held-out participant. The $k = 16$ model of the same curve is the reference, so every paired difference $\Delta_{p,k} = \mathrm{F1}_{p,k} - \mathrm{F1}_{p,16}$ compares models with identical training budgets. Non-inferiority at margin $\delta$ requires the lower bound of the 95\% participant-cluster bootstrap interval of $\bar\Delta_k$ to exceed $-\delta$. We define $k^* = \min\{k < 16 : \mathrm{LCB}_{95}(\bar\Delta_k) > -\delta\}$, recorded as ``none'' if no $k$ qualifies. \rev{The primary margin $\delta = 0.02$ is an absolute difference in class-averaged F1 on a 0--1 scale. It was fixed before the curve was computed as the largest loss we would accept in exchange for fewer electrodes. Relative to the full-array curve reference it is about \MarginPctOfRef\% of performance, and about \MarginPctAboveChance\% of the distance above chance (1/16 under uniform guessing). It is an engineering tolerance: it has not been validated against clinical outcomes or user performance. Because $k^*$ depends strongly on $\delta$, the prespecified sensitivity margins $\delta = 0.01$ and $0.03$ are reported so that readers can apply a different tolerance. The plan also required Holm correction within H1. The primary analysis above did not apply it; we therefore also report, post hoc, one-sided cluster-bootstrap $p$-values for $H_0\colon \bar\Delta_k \le -\delta$ with Holm correction over the 15 tested $k$ at $\alpha = 0.025$ (Section~\ref{sec:deviations}).}

\subsection{Comparison with other selection methods (H2)}
At $k \in \{4, 6, 8, 12\}$ we compared six rankings, each computed on outer-training data only:
\begin{itemize}
\item SR-CFA ($\lambda = 0.5$);
\item cross-fitted SHAP ($\lambda = 0$);
\item permutation importance\rev{: from a 40,000-window stratified subsample of the outer-training data, a model was fitted on 20,000 windows and each of the 160 features was permuted in 8,000 of the remaining windows (3 repeats, Macro-F1 scoring); a channel's score is the mean over its 10 features};
\item mutual information \rev{between each feature and the gesture label, estimated on the same 20,000 windows; a channel's score is the mean over its 10 features};
\item a geometric baseline, which selects electrodes by farthest-point sampling on the $2 \times 8$ montage \rev{(circumferential arc distance on an assumed 24~cm circumference, 2~cm between rings; deterministic start at F1, ties to the lower index). It uses no data};
\item random subsets: 100 draws per participant at each $k$, and 500 at the prespecified $k^*$ \rev{($k$ electrodes without replacement within a draw; draws independent, so a subset can recur; seed fixed per participant)}.
\end{itemize}
The value $\lambda = 0.5$ for this comparison had no recorded rationale. Because the P4 comparison includes $\lambda = 0$ as well, both rankings are reported. Drawing hundreds of random subsets per participant was only tractable with a lighter comparison model (XGBoost with 50 trees, trained on a 150,000-window stratified subsample), applied identically to every method. Its absolute scores are therefore lower than those of the sparse curve and are compared only within the comparison. H2 contrasts each attribution ranking with the participant's mean over its random subsets (Holm over 8 tests). We also report where the SR-CFA subset falls within each participant's random distribution. \rev{Here and below, ``SR-CFA'' refers to $\lambda = 0.5$ and ``cross-fitted SHAP'' to $\lambda = 0$.}

\subsection{Stability regularisation (H3) and faithfulness (H4)}
H3 compared each $\lambda > 0$ with $\lambda = 0$ on the sparse curve at $k \in \{4, 6, 8, 12\}$ (Holm over 12 tests). For H4 we used deletion with retraining~\cite{hooker2019benchmark}. For $n \in \{2, 4, \ldots, 14\}$ \rev{channels retained}, the model was retrained after removing either the $16 - n$ highest-ranked channels (\rev{Remove-Top}: keeping the $n$ lowest) or the $16 - n$ lowest-ranked (\rev{Remove-Bottom}: keeping the top $n$). \rev{Remove-Random} is the mean over the random subsets above, which exist for $n \in \{4, 6, 8, 12\}$. All three arms used the same comparison model and training subsample. As a check on protocol equivalence, the \rev{Remove-Bottom} arm selects the same electrodes as the H2 cross-fitted SHAP subsets in all \FaithProtocolCells\ matching participant-by-$k$ cells. Its Macro-F1 agrees with them to within a mean absolute difference of \FaithProtocolMeanDiff\ (maximum \FaithProtocolMaxDiff). The residual differences are numerical: z-scoring was applied per subset rather than once over all features. A faithful ranking gives \rev{Remove-Top $<$ Remove-Random $<$ Remove-Bottom}. Paired contrasts were Holm-corrected within H4.

\subsection{Statistics}
The participant is the statistical unit ($n = 43$; $n = 10$ for DB5\rev{; $n = 44$ for putEMG. We use ``participant'' for GRABMyo and putEMG and ``subject'' for DB5, following each dataset's documentation}). Windows are never treated as independent. Confidence intervals are 95\% percentile intervals from 10,000 cluster-bootstrap resamples of participants. Paired comparisons use the Wilcoxon signed-rank test with the matched-pairs rank-biserial correlation $r_{rb}$ as effect size. Holm correction is applied within each confirmatory family. Exploratory analyses are uncorrected and labelled as such. \rev{A non-significant difference is reported as ``no statistically detectable difference'', never as equivalence; no equivalence test was performed.} No participant was excluded.

\subsection{Exploratory analyses}
\rev{Exploratory analyses are described in the order in which their results are reported.}

\rev{\subsubsection{Ranking stability}} We measured the agreement of rankings computed from single participants' attributions, from the five inner folds, and from the 43 outer folds. \rev{A participant's profile is the mean of their normalised out-of-fold attribution over the outer folds in which they were an inner-held-out participant. Agreement between two rankings is Spearman's $\rho$ and Kendall's $\tau$ over the 16 channel ranks, and Jaccard@$k$, the size of the intersection of the two top-$k$ sets divided by the size of their union; single-participant agreement uses all 903 participant pairs.} Outer folds share about 98\% of their training participants, so their agreement is only an upper bound.

\rev{\subsubsection{Cross-session}} XGBoost (full configuration) was trained on session~1 of $R_{-p}$ (CS1) and tested on each session of $p$. It was also trained on sessions 1 and 2 (CS2) and tested on session 3 of $p$. CS2 has roughly twice the training data of CS1, so the two conditions differ in both session diversity and quantity.

\rev{\subsubsection{Bipolar montage}} To express the selection unit as a \rev{bipolar pair}, we derived $B_j = F_j - F_{j+8}$ ($j = 1, \ldots, 8$) and re-extracted features. \rev{These are derived pairs of the monopolar recording, not physical bipolar electrodes.} The count-feature threshold was re-derived for this signal ($5 \times 10^{-3}$~mV), because the bipolar signal is about half the monopolar amplitude. \rev{How this value was derived is not fully recorded: on the documented probe records the bipolar median is \LOOBipPooled${}\times 10^{-3}$~mV, for which the stated rule gives $\LOOBipRule \times 10^{-3}$~mV. The leave-one-out medians vary little, so the value is not driven by any single held-out participant, but its provenance is a limitation of this exploratory analysis.} We then repeated the LOSO baseline, the cross-fitted ranking and the sparse curve\rev{. The reference is the $k = 8$ model of the same curve (200,000-window training subsample), not the full bipolar baseline trained on all windows.}

\rev{\subsubsection{NinaPro DB5}} The LOSO baseline, cross-fitted ranking ($K = 5$) and sparse curve were repeated on DB5, using all outer-training windows and the $k = 16$ model as reference. \rev{The plan classed DB5 as exploratory, so its primary decision is uncorrected; the Holm-adjusted result is also reported.}

\rev{\subsubsection{putEMG} The same procedure was applied to putEMG: LOSO over 44 participants; XGBoost with the headline configuration trained on all outer-training windows (at most \PutTrainWindowsMax, below the 200,000-window cap); cross-fitted ranking ($K = 5$, all four $\lambda$); and the $\lambda = 0$ sparse curve for $k = 1, \ldots, 24$ with the $k = 24$ model as reference. Non-inferiority used the same bootstrap, margins and $k^*$ rule, with Holm correction over the 23 tested $k$ also reported. Ten random subsets per participant at $k \in \{6, 12, 18\}$ were compared with the ranked subsets (Wilcoxon, Holm over 3). No hyperparameter or training-set size was calibrated on putEMG. Because the montage has 24 electrodes in three bands, $k$ values are not comparable with GRABMyo's.}

\rev{\subsubsection{Fixed layout}} In a post-hoc analysis, a single 12-electrode layout was derived from a cross-fitted ranking ($\lambda = 0.5$) on the 22 odd-numbered participants. Models with 16 and with 12 electrodes were trained on those participants and tested on the 21 even-numbered participants. This is one split, not a LOSO estimate.

\rev{\subsubsection{Feature redundancy}} IEMG equals 512 times MAV for these windows. We repeated the comparison-model LOSO with IEMG removed and compared both Macro-F1 and channel attribution.

\subsection{Deviations from the analysis plan}
\label{sec:deviations}
The following changes were made after results existed. None alters a prespecified decision rule.
\begin{enumerate}
\item \rev{Holm correction within H1, required by the plan, was not implemented in the primary analysis. It was added post hoc and gives the same $k^*$ at all three margins (Supplementary Table~\ref{tab:h1_holm}). A fixed-sequence analysis (testing $k = 15, 14, \ldots$ and stopping at the first failure) is reported only as a post-hoc sensitivity analysis.}
\item \rev{Training-set sizes were calibrated on held-out folds 1 and 7 before the confirmatory runs.}
\item \rev{The leave-one-out check of the count-feature thresholds was added post hoc.}
\item $\lambda$ was not selected; see above.
\item The H2--H4 families were specified in detail (tests and $k$ sets) after the sparse curve and selection-baseline results had been computed.
\item The comparison model for H2 and H4 was reduced from 150 to 50 trees for compute reasons before any baseline result existed.
\item An earlier faithfulness run mixed arms trained with different budgets and was discarded, and all arms were recomputed under one protocol.
\item An earlier bipolar run used a count threshold inherited from the monopolar signal, and was discarded and recomputed.
\item The fixed-layout analysis was not planned.
\item The window-length sensitivity and the 17-class task, both listed as exploratory in the plan, were not run.
\item \rev{The putEMG replication was not in the plan and was added after all other results existed.}
\end{enumerate}

\section{Results}
\label{sec:results}

\subsection{Full-channel subject-independent performance}
With all 16 electrodes, XGBoost reached a Macro-F1 of \FullXGB\ (95\% CI \FullXGBLCB--\FullXGBUCB) over the 43 held-out participants. The RBF-SVM reached \FullSVM\ (\FullSVMLCB--\FullSVMUCB) and the random forest \FullRF\ (\FullRFLCB--\FullRFUCB) (Table~\ref{tab:full_channel}, Fig.~\ref{fig:full}). \rev{No statistically detectable difference was observed between XGBoost and the SVM}: the paired difference was $\XGBvsSVMDiff$ (95\% CI $\XGBvsSVMLCB$ to $\XGBvsSVMUCB$), $p_{\mathrm{Holm}} \XGBvsSVMP$, $r_{rb} = \XGBvsSVMR$. Both exceeded the random forest ($p_{\mathrm{Holm}} \XGBvsRFP$ and $\SVMvsRFP$; $r_{rb} = \XGBvsRFR$ and $\SVMvsRFR$). The SVM result is a lower bound because of its training subsample (Methods). The spread between participants was far larger than any difference between models: XGBoost ranged from \FullXGBMin\ to \FullXGBMax\ across participants.

\begin{table*}[t]
\centering
\small
\caption{Full 16-channel leave-one-subject-out performance on GRABMyo (43 participants, 16 active gestures; chance Macro-F1 $\approx$ 0.063). CI: participant cluster bootstrap.}
\label{tab:full_channel}
\begin{tabular}{llcccc}
\toprule
Model & Training windows & Macro-F1 & 95\% CI & SD & Range \\
\midrule
XGBoost & all ($\approx$550k) & 0.542 & 0.517--0.567 & 0.085 & 0.385--0.702 \\
RBF-SVM & 40k & 0.538 & 0.513--0.562 & 0.085 & 0.383--0.718 \\
Random forest & 200k & 0.502 & 0.475--0.528 & 0.090 & 0.346--0.670 \\
\bottomrule
\end{tabular}
\par\smallskip\parbox{0.95\linewidth}{\footnotesize Paired comparisons (Wilcoxon signed-rank, Holm over 3): Random forest $-$ RBF-SVM = $-$0.0362 (95\% CI $-$0.0444 to $-$0.0284), $p_{\mathrm{Holm}}$ $<$ 0.001, $r_{rb}$ = -0.97; Random forest $-$ XGBoost = $-$0.0408 (95\% CI $-$0.0475 to $-$0.0340), $p_{\mathrm{Holm}}$ $<$ 0.001, $r_{rb}$ = -0.98; RBF-SVM $-$ XGBoost = $-$0.0045 (95\% CI $-$0.0107 to +0.0015), $p_{\mathrm{Holm}}$ = 0.154, $r_{rb}$ = -0.25.}
\end{table*}


\begin{figure*}[t]
\centering
\includegraphics[width=\textwidth]{figures/fig3_full_channel.pdf}
\caption{Full-channel LOSO performance (43 participants; each point is one held-out participant; grey lines join the same participant; black bar = median; diamond = mean with 95\% cluster-bootstrap CI). (a) Three classifiers. (b) Paired XGBoost $-$ SVM difference per participant, sorted; grey band = 95\% CI of the mean difference.}
\label{fig:full}
\end{figure*}

\subsection{H1: a 12-electrode subset is non-inferior under the prespecified margin}
Channel attribution was strongly concentrated and highly reproducible across outer folds. The same four electrodes (\TopFourSet) formed the top four in \TopFourFolds\ of 43 folds, and they carried \TopFourShare\% of the summed attribution. The eight lowest-ranked electrodes carried \BottomEightShare\% (Fig.~\ref{fig:attribution}). Electrodes \DroppedFourSet\ were the four excluded at $k = 12$ in \DroppedFourFolds\ of 43 folds.

Performance declined smoothly as electrodes were removed (Fig.~\ref{fig:curve}, Table~\ref{tab:sparse_ni}). The $k = 16$ reference of the curve (\CurveKSixteen) differs from the full-data result above because every curve model uses the same fixed 200,000-window training subsample. At $\delta = 0.02$, the smallest non-inferior subset was $k^* = \KStar$ electrodes: $\bar\Delta_{12} = \DeltaKTwelve$ (95\% CI $\DeltaKTwelveLCB$ to $\DeltaKTwelveUCB$). At $k = 11$ the lower bound ($\DeltaKElevenLCB$) crossed the margin. Every $k$ from 12 to 15 was non-inferior, so the result is not an isolated crossing. \rev{The plan specified Holm correction within H1, which the primary analysis did not apply. Applied post hoc (one-sided cluster-bootstrap $p$, 15 values of $k$), it gave the same $k^*$ at all three margins ($\KStarHolmPtOne$, $\KStarHolm$ and $\KStarHolmPtThree$ at $\delta = 0.01$, 0.02 and 0.03; $k = 12$: $p_{\mathrm{Holm}} \HolmPKTwelve$; Supplementary Table~\ref{tab:h1_holm}). A post-hoc fixed-sequence analysis from $k = 15$ downward also gave $k^* = \KStarFixSeq$.} The sensitivity margins gave $k^* = \KStarMarginPtOne$ at $\delta = 0.01$ and $k^* = \KStarMarginPtThree$ at $\delta = 0.03$. Removing four electrodes (\KStarReductionPct\% of the array) is therefore within the prespecified tolerance. Larger reductions are not: halving the array to $k = 8$ cost $\DeltaKEight$ (95\% CI $\DeltaKEightLCB$ to $\DeltaKEightUCB$), and $k = 4$ cost $\DeltaKFour$. All paired differences at $k < 16$ were negative. The claim is non-inferiority within a stated engineering tolerance, not equivalence.

\begin{table*}[t]
\centering
\small
\caption{Sparse sensor curve and non-inferiority (cross-fitted ranking, $\lambda = 0$; 43 participants; every $k$ retrained from scratch on a fixed 200k-window outer-training subsample). $\Delta_k$ = paired difference to the $k = 16$ model of the same curve. Non-inferior when the lower 95\% bound exceeds $-\delta$; $\delta = 0.02$ is primary (bold). \rev{Holm: non-inferiority at $\delta = 0.02$ after the plan's Holm correction over the 15 tested $k$, applied post hoc (Supplementary Table~\ref{tab:h1_holm}).}}
\label{tab:sparse_ni}
\begin{tabular}{rccccccc}
\toprule
$k$ & Macro-F1 (mean $\pm$ SD) & $\Delta_k$ & 95\% CI & $\delta$=0.01 & \textbf{$\delta$=0.02 (primary)} & $\delta$=0.03 & \rev{Holm (0.02)} \\
\midrule
16 & 0.541 $\pm$ 0.084 & reference &  &  &  &  &  \\
15 & 0.538 $\pm$ 0.083 & $-$0.0036 & $-$0.0056 to $-$0.0017 & yes & \textbf{yes} & yes & \rev{yes} \\
14 & 0.535 $\pm$ 0.083 & $-$0.0060 & $-$0.0081 to $-$0.0039 & yes & \textbf{yes} & yes & \rev{yes} \\
13 & 0.530 $\pm$ 0.083 & $-$0.0109 & $-$0.0132 to $-$0.0085 & no & \textbf{yes} & yes & \rev{yes} \\
12 & 0.528 $\pm$ 0.084 & $-$0.0132 & $-$0.0160 to $-$0.0105 & no & \textbf{yes} & yes & \rev{yes} \\
11 & 0.521 $\pm$ 0.084 & $-$0.0202 & $-$0.0245 to $-$0.0160 & no & \textbf{no} & yes & \rev{no} \\
10 & 0.515 $\pm$ 0.083 & $-$0.0258 & $-$0.0311 to $-$0.0205 & no & \textbf{no} & no & \rev{no} \\
9 & 0.510 $\pm$ 0.086 & $-$0.0316 & $-$0.0385 to $-$0.0245 & no & \textbf{no} & no & \rev{no} \\
8 & 0.505 $\pm$ 0.085 & $-$0.0366 & $-$0.0440 to $-$0.0289 & no & \textbf{no} & no & \rev{no} \\
7 & 0.500 $\pm$ 0.085 & $-$0.0412 & $-$0.0488 to $-$0.0331 & no & \textbf{no} & no & \rev{no} \\
6 & 0.480 $\pm$ 0.083 & $-$0.0609 & $-$0.0711 to $-$0.0507 & no & \textbf{no} & no & \rev{no} \\
5 & 0.428 $\pm$ 0.076 & $-$0.1132 & $-$0.1276 to $-$0.0993 & no & \textbf{no} & no & \rev{no} \\
4 & 0.397 $\pm$ 0.071 & $-$0.1444 & $-$0.1590 to $-$0.1306 & no & \textbf{no} & no & \rev{no} \\
3 & 0.340 $\pm$ 0.057 & $-$0.2014 & $-$0.2211 to $-$0.1820 & no & \textbf{no} & no & \rev{no} \\
2 & 0.299 $\pm$ 0.057 & $-$0.2420 & $-$0.2603 to $-$0.2230 & no & \textbf{no} & no & \rev{no} \\
1 & 0.181 $\pm$ 0.040 & $-$0.3602 & $-$0.3831 to $-$0.3368 & no & \textbf{no} & no & \rev{no} \\
\bottomrule
\end{tabular}
\end{table*}


\begin{figure*}[t]
\centering
\includegraphics[width=\textwidth]{figures/fig4_sparse_curve.pdf}
\caption{Sparse sensor curve (43 participants). (a) Mean Macro-F1 of models retrained on the top-$k$ electrodes of each fold's cross-fitted ranking, for the four values of $\lambda$ (curves overlap). Band = 95\% cluster-bootstrap CI for $\lambda = 0$. (b) Paired difference to the $k = 16$ model with 95\% CI. \rev{The thick dark dashed line marks the primary margin ($-\delta$, $\delta = 0.02$); thin grey lines mark the sensitivity margins (0.01, 0.03).}}
\label{fig:curve}
\end{figure*}

\subsection{H2: attribution-ranked subsets \rev{exceed the mean of random subsets}}
Under the shared comparison model (Table~\ref{tab:baselines}, Fig.~\ref{fig:baselines}), SR-CFA subsets outperformed the mean of random subsets at every tested $k$ (Holm-corrected, 8 tests). The advantage was largest at intermediate sizes: $\HTwoKSixDiff$ Macro-F1 at $k = 6$ (95\% CI $\HTwoKSixLCB$ to $\HTwoKSixUCB$, $p_{\mathrm{Holm}} \HTwoKSixP$, $r_{rb} = \HTwoKSixR$) and $\HTwoKEightDiff$ at $k = 8$. It was small at $k = 4$ ($\HTwoKFourDiff$) and at $k = 12$ ($\HTwoKTwelveDiff$, $p_{\mathrm{Holm}} \HTwoKTwelveP$). Within each participant's own random subsets, the SR-CFA subset had a median percentile of \PctKSix\ at $k = 6$, \PctKEight\ at $k = 8$, \PctKFour\ at $k = 4$ and \PctKTwelve\ at $k = 12$. It reached the 95th percentile for only \PctKAboveNinetyFiveSix\ of 43 participants at $k = 6$ and \PctKAboveNinetyFiveTwelve\ at $k = 12$. \rev{Attribution-ranked subsets exceeded the mean of the sampled random subsets at every $k$, but rarely fell in the upper tail ($\geq$ 95th percentile) of each participant's sampled random-subset distribution.}

\rev{In exploratory, uncorrected contrasts, attribution showed limited additional benefit over simpler alternatives. No statistically detectable difference from permutation importance was observed at $k = 4$, 8 or 12} (e.g.\ $k = 8$: $\VSPermKEightDiff$, $p \VSPermKEightP$)\rev{; SR-CFA was higher at $k = 6$} ($\VSPermKSixDiff$, $p \VSPermKSixP$). The geometric baseline, which spreads electrodes evenly around both rings without using any data, \rev{showed no statistically detectable difference from SR-CFA at} $k = 4$ ($\VSGeomKFourDiff$, $p \VSGeomKFourP$), $k = 8$ ($\VSGeomKEightDiff$, $p \VSGeomKEightP$) and $k = 12$ ($\VSGeomKTwelveDiff$, $p \VSGeomKTwelveP$)\rev{; SR-CFA was higher at $k = 6$} ($\VSGeomKSixDiff$, $p \VSGeomKSixP$). \rev{These contrasts do not show that the methods are equivalent.} Mutual information was the weakest ranking at every $k \le 8$.

\begin{table*}[t]
\centering
\small
\caption{Channel-selection methods at matched $k$ (43 participants). All methods use one comparison model (XGBoost, 50 trees, 150k-window outer-training subsample), so values are comparable within the table but lower than Table~\ref{tab:sparse_ni}. H2 $p$ values Holm-corrected over 8 tests; $^\dagger$exploratory, uncorrected.}
\label{tab:baselines}
\begin{tabular}{lcccc}
\toprule
Method & $k$ = 4 & $k$ = 6 & $k$ = 8 & $k$ = 12 \\
\midrule
SR-CFA ($\lambda$=0.5) & 0.383 $\pm$ 0.071 & 0.453 $\pm$ 0.080 & 0.469 $\pm$ 0.081 & 0.485 $\pm$ 0.079 \\
Cross-fitted SHAP ($\lambda$=0) & 0.382 $\pm$ 0.071 & 0.453 $\pm$ 0.080 & 0.468 $\pm$ 0.081 & 0.485 $\pm$ 0.079 \\
Permutation importance & 0.385 $\pm$ 0.070 & 0.440 $\pm$ 0.080 & 0.466 $\pm$ 0.080 & 0.485 $\pm$ 0.079 \\
Geometric farthest-point & 0.383 $\pm$ 0.069 & 0.419 $\pm$ 0.068 & 0.465 $\pm$ 0.072 & 0.478 $\pm$ 0.076 \\
Mutual information & 0.280 $\pm$ 0.049 & 0.360 $\pm$ 0.066 & 0.396 $\pm$ 0.068 & 0.480 $\pm$ 0.079 \\
Random subsets (mean) & 0.366 $\pm$ 0.059 & 0.416 $\pm$ 0.068 & 0.445 $\pm$ 0.071 & 0.478 $\pm$ 0.077 \\
\midrule SR-CFA $-$ random (H2) & +0.0172, \rev{$p$ = 0.007} & +0.0370, \rev{$p$ $<$ 0.001} & +0.0243, \rev{$p$ $<$ 0.001} & +0.0074, \rev{$p$ $<$ 0.001} \\
SR-CFA $-$ perm$^\dagger$ & $-$0.0018, \rev{$p$ = 0.678} & +0.0126, \rev{$p$ = 0.009} & +0.0033, \rev{$p$ = 0.069} & +0.0001, \rev{$p$ = 0.995} \\
SR-CFA $-$ geom$^\dagger$ & $-$0.0005, \rev{$p$ = 0.986} & +0.0341, \rev{$p$ $<$ 0.001} & +0.0047, \rev{$p$ = 0.267} & +0.0070, \rev{$p$ = 0.058} \\
SR-CFA percentile in random draws (median) & 65 (100 draws) & 89 (100 draws) & 82 (100 draws) & 66 (500 draws) \\
\bottomrule
\end{tabular}
\end{table*}


\begin{figure*}[t]
\centering
\includegraphics[width=\textwidth]{figures/fig5_selection_baselines.pdf}
\caption{Selection methods at matched $k$ (43 participants; comparison model shared by all methods). (a) Mean Macro-F1 with 95\% cluster-bootstrap CI. (b) Percentile of the SR-CFA subset within each participant's own random subsets (one point per participant; bar = median; 100 draws, 500 at $k = 12$).}
\label{fig:baselines}
\end{figure*}

\subsection{H3: the stability penalty had no measurable effect}
The instability term was small relative to attribution magnitude: the median $U_c/A_c$ was \UOverAMedian\% (maximum \UOverAMax\%). Penalising it therefore rarely changed a ranking. Against $\lambda = 0$, the top-12 set changed for \TopTwelveDiffersLambdaPtTwoFive, \TopTwelveDiffersLambdaPtFive\ and \TopTwelveDiffersLambdaOne\ of 43 folds at $\lambda = 0.25$, 0.5 and 1. Across all four $k$ tested, at least \HThreeMinTied\ of 43 participants received the identical subset under both rankings. None of the \HThreeNTests\ contrasts was significant after Holm correction (largest mean difference \HThreeMaxAbsDiff; smallest $p_{\mathrm{Holm}} \HThreeMinP$; Supplementary Table~\ref{tab:lambda}). The four curves in Fig.~\ref{fig:curve}a are indistinguishable (maximum spread between $\lambda$ values \LambdaCurveMaxSpread), and $k^*$ was identical under every $\lambda$ and margin. H3 is not supported. \rev{Under this protocol the stability penalty produced no measurable improvement, consistent with the small instability term.}

\subsection{H4: \rev{deletion with retraining supports faithfulness}}
When all three deletion arms were trained under one protocol, the expected ordering held at every retained size (Table~\ref{tab:faithfulness}, Fig.~\ref{fig:faithfulness}). Deleting the highest-ranked electrodes hurt more than deleting random ones, and random deletion hurt more than deleting the lowest-ranked. \FaithNSig\ of the \FaithNTests\ Holm-corrected contrasts were significant. \rev{With six channels retained, Remove-Top} scored \FaithTopKSix, \rev{Remove-Random} \FaithRandomKSix\ and \rev{Remove-Bottom} \FaithBottomKSix\ (\rev{Remove-Top $-$ Remove-Random} $\FaithTRKSixDiff$, $p_{\mathrm{Holm}} \FaithTRKSixP$; \rev{Remove-Random $-$ Remove-Bottom} $\FaithRBKSixDiff$, $p_{\mathrm{Holm}} \FaithRBKSixP$). The only non-significant contrast was \rev{Remove-Top versus Remove-Bottom with 14 channels retained} ($\FaithTBKFourteenDiff$, $p_{\mathrm{Holm}} \FaithTBKFourteenP$), where only two electrodes are removed. \rev{The separation between deletion orders was largest at intermediate sizes and small at the extremes.}

\begin{figure}[t]
\centering
\includegraphics[width=\columnwidth]{figures/fig7_faithfulness.pdf}
\caption{Deletion-with-retraining test of attribution faithfulness (43 participants; mean and 95\% cluster-bootstrap CI). All arms use the same comparison model and training subsample. \rev{Remove-Random} is available where random subsets were sampled ($n$ = 4, 6, 8, 12).}
\label{fig:faithfulness}
\end{figure}

\begin{table*}[t]
\centering
\small
\caption{Deletion (remove-and-retrain) faithfulness test (H4; 43 participants). Every arm uses the same comparison model and training subsample; \rev{Remove-Random} averages 100 (500 at $k$ = 12) draws per participant and exists only where those draws were sampled. Differences are paired means with Holm-corrected $p$ in parentheses.}
\label{tab:faithfulness}
\begin{tabular}{rcccccc}
\toprule
\rev{Retained} & \rev{Remove-Top} & \rev{Remove-Random} & \rev{Remove-Bottom} & Top$-$Bottom & Top$-$Random & Random$-$Bottom \\
\midrule
2 & 0.263 & -- & 0.292 & $-$0.0297 ($<$0.001) & -- & -- \\
4 & 0.351 & 0.366 & 0.382 & $-$0.0314 ($<$0.001) & $-$0.0151 (0.010) & $-$0.0163 (0.013) \\
6 & 0.376 & 0.416 & 0.454 & $-$0.0775 ($<$0.001) & $-$0.0398 ($<$0.001) & $-$0.0377 ($<$0.001) \\
8 & 0.408 & 0.445 & 0.469 & $-$0.0618 ($<$0.001) & $-$0.0377 ($<$0.001) & $-$0.0242 ($<$0.001) \\
10 & 0.445 & -- & 0.475 & $-$0.0297 ($<$0.001) & -- & -- \\
12 & 0.455 & 0.478 & 0.485 & $-$0.0302 ($<$0.001) & $-$0.0228 ($<$0.001) & $-$0.0074 ($<$0.001) \\
14 & 0.487 & -- & 0.491 & $-$0.0038 (0.150) & -- & -- \\
\bottomrule
\end{tabular}
\end{table*}


\subsection{Stability of the ranking across people}
Rankings computed from individual participants' attributions agreed closely with each other (Table~\ref{tab:stability}): mean Spearman $\rho$ = \StabSubjRho\ $\pm$ \StabSubjRhoSD\ over \StabSubjPairs\ pairs, and Jaccard overlap of the top six = \StabSubjJacSix. Rankings from disjoint inner folds agreed at $\rho$ = \StabInnerRho. The outer-fold agreement ($\rho$ = \StabOuterRho) is an upper bound only, because those training sets share about 98\% of their participants. Across folds, the highest- and lowest-ranked electrodes kept their positions, and the rank ranges were widest in the middle of the ordering (Fig.~\ref{fig:attribution}b).

\begin{table*}[t]
\centering
\small
\caption{Agreement of channel rankings ($\lambda = 0$). Values are means over pairs.}
\label{tab:stability}
\begin{tabular}{lcccccc}
\toprule
Comparison & Pairs & Spearman $\rho$ & Kendall $\tau$ & Jaccard@4 & Jaccard@6 & Jaccard@8 \\
\midrule
Per-subject attribution (between people) & 903 & 0.927 $\pm$ 0.038 & 0.801 & 0.752 & 0.905 & 0.803 \\
Inner folds (disjoint subject groups) & 430 & 0.948 $\pm$ 0.021 & 0.840 & 0.816 & 0.992 & 0.838 \\
Outer folds (upper bound$^*$) & 903 & 0.987 $\pm$ 0.008 & 0.943 & 0.888 & 1.000 & 0.923 \\
\bottomrule
\end{tabular}
\par\smallskip\parbox{0.95\linewidth}{\footnotesize $^*$Outer training sets share about 98\% of their participants, so their agreement is an upper bound, not an estimate of between-subject stability.}
\end{table*}


\begin{figure*}[t]
\centering
\includegraphics[width=\textwidth]{figures/fig6_attribution.pdf}
\caption{Cross-fitted channel attribution ($\lambda = 0$, 43 outer folds). (a) Mean share of attribution per electrode, drawn on the $2 \times 8$ montage (ring~2 electrodes are aligned with, and 2~cm from, ring~1 electrodes at the same position). (b) Rank of each electrode across folds (median and range).}
\label{fig:attribution}
\end{figure*}

\subsection{Exploratory analyses}

\begin{figure*}[t]
\centering
\includegraphics[width=\textwidth]{figures/fig8_generalisation.pdf}
\caption{\rev{Exploratory analyses. (a) Cross-session: grey lines join the same held-out participant across conditions (bar = median); all conditions are also subject-independent. (b--d) Paired difference to the full-array reference of the same curve, $\bar\Delta_k$ with 95\% cluster-bootstrap CI, for (b) NinaPro DB5 (10 subjects, reference $k = 16$), (c) GRABMyo derived bipolar pairs (43 participants, $\lambda = 0$, reference $k = 8$) and (d) putEMG (44 participants, reference $k = 24$). Thick dashed line = primary margin $-0.02$; filled markers = lower bound above the margin.}}
\label{fig:generalisation}
\end{figure*}

\paragraph{Cross-session}
Training on session~1 and testing on sessions 1, 2 and 3 gave \CSOneSOne, \CSOneSTwo\ and \CSOneSThree\ (Table~\ref{tab:cross_session}, Fig.~\ref{fig:generalisation}a). \rev{Mean Macro-F1 on sessions 2 and 3 was numerically higher than on session 1. The session-2 difference was nominally significant} ($\CSSTwoVsSOneDiff$, 95\% CI $\CSSTwoVsSOneLCB$ to $\CSSTwoVsSOneUCB$, uncorrected $p \CSSTwoVsSOneP$), and \rev{no statistically detectable difference was observed for session 3} ($\CSSThreeVsSOneDiff$, $p \CSSThreeVsSOneP$). Because every condition is also subject-independent, this does not show that sessions are interchangeable. It shows that, once the cross-user gap is present, no additional cross-day loss was detectable. Training on sessions 1 and 2 improved session-3 performance over training on session 1 alone ($\CSTwoVsOneDiff$, 95\% CI $\CSTwoVsOneLCB$ to $\CSTwoVsOneUCB$, $p \CSTwoVsOneP$). This design cannot separate day diversity from the doubled amount of training data.

\paragraph{Bipolar montage}
With the eight derived bipolar pairs as the selection unit, full-array LOSO Macro-F1 was \BipFull\ $\pm$ \BipFullSD, lower than the monopolar array (Fig.~\ref{fig:generalisation}c, \rev{Supplementary Table~\ref{tab:bipolar}}). No pair could be removed within the primary margin. With seven pairs, $\bar\Delta_7 = \BipDeltaKSeven$, and its lower bound ($\BipDeltaKSevenLCB$) fell just outside $-0.02$. It did so under every one of 200 alternative bootstrap seeds (range $\BipSeedLCBMinSeven$ to $\BipSeedLCBMaxSeven$). $k^*$ was therefore \BipKStarMarginPtTwo\ at $\delta = 0.02$, \BipKStarMarginPtOne\ at 0.01, and \BipKStarMarginPtThree\ pairs at 0.03, identically for all $\lambda$. The ranking was again stable: pairs \BipTopThreeSet\ formed the top three in \BipTopThreeFolds\ of 43 folds. \rev{With derived bipolar pairs as the unit, no reduction met the primary margin. The reference here is the $k = 8$ model of the bipolar curve (\BipCurveRefK), trained on the same subsample as the other curve points.}

\paragraph{NinaPro DB5}
On DB5 the full 16-channel LOSO Macro-F1 was \DBFiveFull\ $\pm$ \DBFiveFullSD\ (12 movements; chance about 0.083), far below GRABMyo and with wide between-subject spread (\DBFiveFullMin--\DBFiveFullMax) (Fig.~\ref{fig:generalisation}b, Supplementary Table~\ref{tab:replication}). Under the frozen rule, $k^* = \DBFiveKStarMarginPtTwo$ at $\delta = 0.02$ ($\bar\Delta_{12} = \DBFiveDeltaKTwelve$, lower bound $\rev{\DBFiveDeltaKTwelveLCBSix}$). This result is fragile. The lower bound lies on the margin, and $k = 12$ was non-inferior under \DBFiveSeedFrac\% of \DBFiveSeedN\ alternative bootstrap seeds (lower bound range $\DBFiveSeedLCBMin$ to $\DBFiveSeedLCBMax$). Moreover, $k = 13$, 14 and 15 all failed (lower bounds $\DBFiveDeltaKThirteenLCB$, $\DBFiveDeltaKFourteenLCB$, $\DBFiveDeltaKFifteenLCB$). \rev{With Holm correction over the 15 tested $k$ (not required for this exploratory analysis), no subset was non-inferior at $\delta = 0.02$ ($k^* = \DBFiveKStarHolmPtTwo$), and a post-hoc fixed-sequence analysis from $k = 15$ would stop at the first step.} With ten subjects, DB5 does not support a claim that a quarter of the channels can be removed. \rev{The procedure ran on DB5 without modification; the mean curve increased broadly with $k$ but was not monotone.}

\rev{\paragraph{putEMG (exploratory, added post hoc)} With all 24 electrodes, LOSO Macro-F1 over the 44 participants was \PutFull\ (95\% CI \PutFullLCB--\PutFullUCB; range \PutFullMin--\PutFullMax; 7 gestures, chance about 0.143). At $\delta = 0.02$ the smallest non-inferior subset was $k^* = \PutKStarPtTwo$ of 24 electrodes ($\bar\Delta = \PutDeltaKFifteen$, 95\% CI $\PutDeltaKFifteenLCB$ to $\PutDeltaKFifteenUCB$), a reduction of \PutKStarReductionPct\%; at $k = \PutKStarMinusOne$ the lower bound was $\PutDeltaKFourteenLCB$. The passing $k$ values were contiguous up to 23. With Holm correction over the 23 tested $k$, $k^* = \PutKStarHolmPtTwo$. Under 200 alternative bootstrap seeds, $k^*$ was non-inferior in \PutSeedFrac\% of cases. Under the sensitivity margins, $k^*$ was \PutKStarPtOne\ at $\delta = 0.01$ and \PutKStarPtThree\ at $\delta = 0.03$ (Fig.~\ref{fig:generalisation}d). The mean curve increased broadly with $k$ but was not monotone. Ranked subsets exceeded the mean of ten random subsets at all three sizes tested ($k = 6$: $\PutVsRandKSixDiff$ (95\% CI $\PutVsRandKSixLCB$ to $\PutVsRandKSixUCB$), $p_{\mathrm{Holm}} \PutVsRandKSixP$; $k = 12$: $\PutVsRandKTwelveDiff$ (95\% CI $\PutVsRandKTwelveLCB$ to $\PutVsRandKTwelveUCB$), $p_{\mathrm{Holm}} \PutVsRandKTwelveP$; $k = 18$: $\PutVsRandKEighteenDiff$ (95\% CI $\PutVsRandKEighteenLCB$ to $\PutVsRandKEighteenUCB$), $p_{\mathrm{Holm}} \PutVsRandKEighteenP$). The modal top four electrodes were \PutTopFourSet\ (\PutTopFourFolds\ of 44 folds); of the top-8 positions across folds, \PutTopEightElbowPct\%, \PutTopEightMiddlePct\% and \PutTopEightWristPct\% fell in the electrode bands 1--8, 9--16 and 17--24. The top-12 set was identical under all four $\lambda$ in \PutTopTwelveLambdaSame\ of 44 folds. Because putEMG has a different montage, electrode count and gesture set, only the behaviour of the procedure, not the retained fraction, is comparable with GRABMyo.}

\paragraph{Fixed layout (post hoc)}
One 12-electrode layout (removing electrodes \FixedDropped) was derived on \FixedDev\ participants and applied unchanged to \FixedTest\ others. It scored $\FixedDiff$ against the 16-electrode model trained on the same participants (95\% CI $\FixedLCB$ to $\FixedUCB$). This lies within the margin, but it is one post-hoc split and is reported only as a consistency check on the LOSO result.

\paragraph{Feature redundancy}
Removing IEMG, an exact multiple of MAV, changed Macro-F1 by $\RedundDiff$ (95\% CI $\RedundLCB$ to $\RedundUCB$; 43 folds, comparison model). Channel rankings were almost unchanged: the median per-fold Spearman $\rho$ between rankings with and without IEMG was \RedundRho, and the top-12 set was identical in \RedundTopTwelveSame\ of 43 folds. \rev{Removing the duplicated feature therefore had little measurable effect on the channel ranking.}

\section{Discussion}
\label{sec:discussion}

\paragraph{What the prespecified tests show}
Under a margin fixed before the curve was computed, four of the 16 GRABMyo electrodes could be removed while remaining non-inferior to the full array. Five could not. Every $k$ from 12 to 15 met the margin, so the result does not rest on a single \rev{isolated} crossing\rev{, and the plan's Holm correction, applied post hoc, left $k^*$ unchanged}. The margin is an engineering tolerance, not a clinical threshold. Every reduction cost some Macro-F1, and halving the array cost about \DeltaKEight. The practical reading is modest: in a subject-independent setting, this 16-electrode \rev{monopolar} array contains only limited redundancy for 16-gesture recognition. User-dependent studies that retain far fewer channels answer a different question~\cite{phinyomark2018emg,qu2021reduce}, and their results should not be taken as evidence that a \rev{layout intended for many users} can be equally sparse.

Two of the three hypotheses about attribution were supported. Attribution-ranked subsets \rev{exceeded the mean of the random subsets} at every size tested (H2), and deletion with retraining showed the expected ordering at every size (H4). \rev{Deletion with retraining therefore supports that the ranking reflects channels on which the retrained models depend.} The stability penalty, however, \rev{produced no measurable benefit} (H3). With cross-fitting over roughly 34 participants per inner-training set, inner-fold attributions already agreed so closely that the dispersion term was a few percent of the magnitude term. \rev{It changed the top-12 set in at most one fold.} We report SR-CFA as a tested and unsupported variant. \rev{The evidence therefore favours plain cross-fitted TreeSHAP magnitude over the regularised variant;} a related robustness-weighted score has been proposed independently~\cite{xiang2026roshap}.

\paragraph{\rev{Limited additional benefit over simple alternatives}}
The exploratory comparisons temper H2. \rev{Exceeding the mean of random subsets} is a low bar. \rev{In exploratory comparisons, neither a data-free geometric rule nor permutation importance differed detectably from the attribution ranking at 4, 8 and 12 electrodes; attribution was higher at $k = 6$. Absence of a detectable difference is not evidence of equivalence.} Within each participant's own random subsets, the attribution subset was typically above the median but rarely near the top. \rev{One hypothesis consistent with this pattern is that even circumferential coverage accounts for much of what a ranking achieves. Whether a data-free, evenly spread layout is itself non-inferior to the full array was not tested and is a question for future work.}

\paragraph{Interpretation of the ranked electrodes}
The four highest-ranked electrodes (\TopFourSet) were the same in \TopFourFolds\ of 43 folds and in single-participant rankings (mean Spearman $\rho$ = \StabSubjRho). We deliberately do not map them to named muscles. The dataset documents ring positions, not the anatomical placement of each circumferential index\rev{, and no imaging-based registration of electrodes to muscles was available}. Surface electrodes also record compound activity from several muscles at each site~\cite{farina2004extraction}. The robust statement is positional. High attribution concentrated at a few circumferential positions (Fig.~\ref{fig:attribution}a). The top and bottom of the ranking were stable, and the ordering varied mainly among mid-ranked electrodes. \rev{Physiologically, it is plausible that a few circumferential sites capture activity shared by several gestures, but these data cannot attribute that activity to specific muscles.}

\paragraph{Cross-session and replication}
Cross-session results were exploratory and should be read narrowly. \rev{No additional loss on later days was detected on top of the cross-user gap.} This \rev{does not show} that electrode shift is harmless~\cite{hargrove2006effect,kaufmann2010fluctuating,vidovic2016improving}. The gain from training on two sessions is confounded with doubling the training data. On DB5 the whole procedure ran without modification, but \rev{with 10 subjects the estimates were imprecise, and the $k = 12$ decision depended on the bootstrap seed and did not survive Holm correction, with neighbouring $k$ failing. The result neither replicates nor contradicts the GRABMyo reduction.} \rev{On putEMG, a larger cohort with a three-band, 24-electrode montage added after the main analysis, the same procedure also found a non-inferior reduction ($k^* = \PutKStarPtTwo$ of 24). Because this analysis was post hoc and the montage differs, it supports the transferability of the procedure rather than any particular retained fraction.}

\paragraph{Limitations}
The primary unit is a monopolar electrode. A conventional device measures bipolar pairs, and removing four electrodes does not remove four sensors of such a device. The exploratory bipolar re-analysis points the other way: with pairs as the unit, not even one of eight could be removed within the primary margin, although the shortfall at seven pairs was small. The 25\% reduction should therefore be read as a statement about monopolar electrodes on this array, not as evidence that a bipolar armband can shed a quarter of its sensors. The bipolar reference is itself a derived montage, and physical bipolar electrodes may behave differently. All participants were able-bodied, and all analysis was offline\rev{: filtering was zero-phase (non-causal) and windows were 250~ms, so the results are an offline benchmark and do not establish performance in real-time, causal control. Fatigue, limb posture and electrode shift were not isolated. Rankings were computed from full-array models and could shift if models were trained only on a reduced array. No imaging-based registration of electrodes to muscles was available.} The SVM was trained on a subsample and is \rev{likely} a lower bound. The H2 and H4 comparisons used a lighter comparison model than the headline curve. The count-feature threshold was a single scalar set \rev{once rather than per fold (a leave-one-out check found no effect), training-set sizes were calibrated on two held-out folds, and the bipolar threshold's derivation is not fully recorded}. $\lambda$ was not selected, and the detailed H2--H4 test families were specified after the underlying results were computed (Section~\ref{sec:deviations}). \rev{The fixed-layout analysis and the putEMG replication were post hoc, and DB5 had only 10 subjects.} Cross-model attribution, window-length sensitivity and the 17-class task were listed as exploratory in the plan and were not run.

\section{Conclusion}
\label{sec:conclusion}

\rev{We evaluated explainability-guided electrode reduction for subject-independent sEMG gesture recognition using cross-fitted attribution and prespecified primary decision rules.}

\rev{\textit{Supported}, for subject-independent, offline recognition with a two-ring monopolar forearm montage in able-bodied participants: removal of four of 16 electrodes within a prespecified 0.02 Macro-F1 tolerance, unchanged under Holm correction, but not of five; attribution-ranked subsets above the mean of random subsets of the same size; deletion-based support for the faithfulness of the ranking; and no measurable benefit from the stability penalty. In an exploratory, post-hoc replication on putEMG, $k^* = \PutKStarPtTwo$ of 24 electrodes met the same tolerance.}

\rev{\textit{Not established}: clinical performance; reduction of physical bipolar hardware (with derived bipolar pairs, no pair could be removed within the same tolerance); equivalence to, or superiority over, a data-free geometric selection; real-time use; muscle-level localisation; and generalisation beyond these cohorts and montages.}

\rev{Validation in target clinical populations (e.g.\ people with upper-limb loss) and in online, causal control settings is required.}


\section*{Declaration of Competing Interest}
The authors declare that they have no known competing financial interests or personal relationships that could have appeared to influence the work reported in this paper.

\section*{Data and Code Availability}
GRABMyo v1.1.0 is available from PhysioNet (\url{https://physionet.org/content/grabmyo/1.1.0/}), NinaPro DB5 from the NinaPro repository\rev{, and putEMG from its authors' repository~\cite{kaczmarek2019putemg} under a CC BY-NC 4.0 licence}. The frozen analysis plan, decision log, analysis code and every result file cited in this paper are included in the accompanying repository. All numbers in the text and tables are generated from those files.

\appendix
\section{Supplementary tables}
\begin{table}[t]
\centering
\small
\caption{Stability regularisation versus attribution magnitude alone (H3): $Q(\lambda)$ ranking minus $Q(0)$ ranking on the sparse curve. Tied = participants for whom both rankings select the identical subset. Holm over 12 tests.}
\label{tab:lambda}
\begin{tabular}{rrcccc}
\toprule
$\lambda$ & $k$ & Mean diff & 95\% CI & Tied & $p_{\mathrm{Holm}}$ \\
\midrule
0.25 & 4 & $-$0.0002 & $-$0.0005 to +0.0000 & 41/43 & 1.000 \\
0.25 & 6 & +0.0000 & +0.0000 to +0.0000 & 43/43 & 1.000 \\
0.25 & 8 & +0.0001 & +0.0000 to +0.0002 & 42/43 & 1.000 \\
0.25 & 12 & +0.0000 & +0.0000 to +0.0000 & 43/43 & 1.000 \\
0.5 & 4 & +0.0005 & $-$0.0004 to +0.0019 & 39/43 & 1.000 \\
0.5 & 6 & $-$0.0001 & $-$0.0002 to +0.0000 & 42/43 & 1.000 \\
0.5 & 8 & +0.0008 & +0.0000 to +0.0020 & 40/43 & 1.000 \\
0.5 & 12 & $-$0.0000 & $-$0.0001 to +0.0000 & 42/43 & 1.000 \\
1 & 4 & +0.0010 & $-$0.0002 to +0.0028 & 37/43 & 1.000 \\
1 & 6 & $-$0.0001 & $-$0.0002 to +0.0000 & 42/43 & 1.000 \\
1 & 8 & +0.0015 & +0.0004 to +0.0029 & 36/43 & 0.216 \\
1 & 12 & $-$0.0005 & $-$0.0014 to +0.0000 & 41/43 & 1.000 \\
\bottomrule
\end{tabular}
\end{table}

\begin{table*}[t]
\centering
\small
\caption{Cross-session evaluation (exploratory). All conditions are also subject-independent: the test participant is never in training.}
\label{tab:cross_session}
\begin{tabular}{lcccc}
\toprule
Design & Train sessions & Test session & Train windows & Macro-F1 \\
\midrule
CS1 & 1 & 1 & 183{,}456 & 0.492 $\pm$ 0.101 \\
CS1 & 1 & 2 & 183{,}456 & 0.528 $\pm$ 0.088 \\
CS1 & 1 & 3 & 183{,}456 & 0.510 $\pm$ 0.090 \\
CS2 & 1+2 & 3 & 366{,}912 & 0.530 $\pm$ 0.095 \\
\bottomrule
\end{tabular}
\par\smallskip\parbox{0.95\linewidth}{\footnotesize Paired contrasts (exploratory, uncorrected): CS1: test S2 $-$ test S1: +0.0355 (95\% CI +0.0091 to +0.0621), $p$ = 0.024; CS1: test S3 $-$ test S1: +0.0172 (95\% CI $-$0.0038 to +0.0391), $p$ = 0.272; S3 test: CS2 (train S1+S2) $-$ CS1 (train S1): +0.0210 (95\% CI +0.0146 to +0.0274), $p$ $<$ 0.001.}
\end{table*}

\begin{table}[t]
\centering
\small
\caption{\rev{NinaPro DB5 sparse curve (exploratory; cross-fitted ranking, $\lambda = 0$; 10 subjects; reference $k = 16$, mean Macro-F1 0.281; margin $\delta = 0.02$). LCB to 6 decimals because the $k = 12$ decision lies at the margin. NI = lower bound above $-\delta$ (primary, uncorrected); Holm = after Holm correction over 15 $k$; seeds = share of 200 alternative bootstrap seeds under which $k$ is non-inferior.}}
\label{tab:replication}
\begin{tabular}{rcccccc}
\toprule
$k$ & Macro-F1 & $\Delta_k$ & LCB & NI & Holm & Seeds NI \\
\midrule
15 & 0.270 & $-$0.0108 & \rev{$-$0.023773} & no & \rev{no} & 0\% \\
14 & 0.269 & $-$0.0119 & \rev{$-$0.027438} & no & \rev{no} & 0\% \\
13 & 0.266 & $-$0.0143 & \rev{$-$0.031186} & no & \rev{no} & 0\% \\
12 & 0.268 & $-$0.0126 & \rev{$-$0.019994} & yes & \rev{no} & 96\% \\
11 & 0.260 & $-$0.0211 & \rev{$-$0.034121} & no & \rev{no} & 0\% \\
10 & 0.265 & $-$0.0160 & \rev{$-$0.031368} & no & \rev{no} & 0\% \\
9 & 0.260 & $-$0.0202 & \rev{$-$0.036943} & no & \rev{no} & 0\% \\
8 & 0.252 & $-$0.0288 & \rev{$-$0.055369} & no & \rev{no} & 0\% \\
7 & 0.232 & $-$0.0486 & \rev{$-$0.075346} & no & \rev{no} & 0\% \\
6 & 0.231 & $-$0.0498 & \rev{$-$0.081816} & no & \rev{no} & 0\% \\
5 & 0.200 & $-$0.0803 & \rev{$-$0.118422} & no & \rev{no} & 0\% \\
4 & 0.180 & $-$0.1006 & \rev{$-$0.150266} & no & \rev{no} & 0\% \\
3 & 0.162 & $-$0.1185 & \rev{$-$0.174145} & no & \rev{no} & 0\% \\
2 & 0.138 & $-$0.1430 & \rev{$-$0.195307} & no & \rev{no} & 0\% \\
1 & 0.105 & $-$0.1754 & \rev{$-$0.229548} & no & \rev{no} & 0\% \\
\bottomrule
\end{tabular}
\end{table}

\begin{table}[t]
\centering
\small
\caption{\rev{GRABMyo derived bipolar pairs $F_j - F_{j+8}$ (exploratory; 43 participants; $\lambda = 0$). Reference = $k = 8$ model of the same curve (200k-window subsample), mean Macro-F1 0.479; margin $\delta = 0.02$.}}
\label{tab:bipolar}
\begin{tabular}{rccccc}
\toprule
$k$ & Macro-F1 & $\Delta_k$ & 95\% CI & NI & Seeds NI \\
\midrule
7 & 0.463 & $-$0.0155 & $-$0.0203 to $-$0.0105 & no & 0\% \\
6 & 0.441 & $-$0.0376 & $-$0.0456 to $-$0.0301 & no & 0\% \\
5 & 0.405 & $-$0.0737 & $-$0.0841 to $-$0.0639 & no & 0\% \\
4 & 0.354 & $-$0.1245 & $-$0.1395 to $-$0.1101 & no & 0\% \\
3 & 0.309 & $-$0.1703 & $-$0.1857 to $-$0.1553 & no & 0\% \\
2 & 0.239 & $-$0.2394 & $-$0.2608 to $-$0.2171 & no & 0\% \\
1 & 0.174 & $-$0.3047 & $-$0.3295 to $-$0.2783 & no & 0\% \\
\bottomrule
\end{tabular}
\end{table}

\begin{table}[t]
\centering
\small
\caption{\rev{H1 with the plan's Holm correction (applied post hoc; 43 participants, $\lambda = 0$, $\delta = 0.02$). $p$ = one-sided cluster-bootstrap $p$ for $H_0\colon \bar\Delta_k \le -\delta$ (10,000 resamples); Holm over the 15 tested $k$ at $\alpha = 0.025$. Fixed sequence (post-hoc sensitivity): $k = 15, 14, \ldots$ tested in order, stopping at the first failure.}}
\label{tab:h1_holm}
\begin{tabular}{rccccccc}
\toprule
$k$ & $\bar\Delta_k$ & LCB & $p$ & $p_{\mathrm{Holm}}$ & NI & NI Holm & Fixed seq. \\
\midrule
15 & $-$0.0036 & $-$0.0056 & $<$0.001 & 0.002 & yes & yes & yes \\
14 & $-$0.0060 & $-$0.0081 & $<$0.001 & 0.002 & yes & yes & yes \\
13 & $-$0.0109 & $-$0.0132 & $<$0.001 & 0.002 & yes & yes & yes \\
12 & $-$0.0132 & $-$0.0160 & $<$0.001 & 0.002 & yes & yes & yes \\
11 & $-$0.0202 & $-$0.0245 & 0.544 & 1.000 & no & no & no \\
10 & $-$0.0258 & $-$0.0311 & 0.984 & 1.000 & no & no & no \\
9 & $-$0.0316 & $-$0.0385 & 0.999 & 1.000 & no & no & no \\
8 & $-$0.0366 & $-$0.0440 & 1.000 & 1.000 & no & no & no \\
7 & $-$0.0412 & $-$0.0488 & 1.000 & 1.000 & no & no & no \\
6 & $-$0.0609 & $-$0.0711 & 1.000 & 1.000 & no & no & no \\
5 & $-$0.1132 & $-$0.1276 & 1.000 & 1.000 & no & no & no \\
4 & $-$0.1444 & $-$0.1590 & 1.000 & 1.000 & no & no & no \\
3 & $-$0.2014 & $-$0.2211 & 1.000 & 1.000 & no & no & no \\
2 & $-$0.2420 & $-$0.2603 & 1.000 & 1.000 & no & no & no \\
1 & $-$0.3602 & $-$0.3831 & 1.000 & 1.000 & no & no & no \\
\bottomrule
\end{tabular}
\par\smallskip\parbox{0.95\linewidth}{\footnotesize \rev{$k^*$ (unadjusted / Holm / fixed sequence) at $\lambda = 0$: $\delta$ = 0.01: 14 / 14 / 14; $\delta$ = 0.02: 12 / 12 / 12; $\delta$ = 0.03: 11 / 11 / 11; identical for all four $\lambda$.}}
\end{table}

\begin{table}[t]
\centering
\small
\caption{\rev{putEMG sparse curve (exploratory, added post hoc; 44 participants, 24 monopolar electrodes, 7 gestures; cross-fitted ranking, $\lambda = 0$; reference $k = 24$, mean Macro-F1 0.680; margin $\delta = 0.02$). NI = lower bound above $-\delta$; Holm over 23 $k$; seeds = share of 200 alternative bootstrap seeds.}}
\label{tab:putemg}
\begin{tabular}{rcccccc}
\toprule
$k$ & Macro-F1 & $\bar\Delta_k$ & 95\% CI & NI & Holm & Seeds NI \\
\midrule
23 & 0.680 $\pm$ 0.118 & $-$0.0000 & $-$0.0022 to +0.0020 & yes & yes & 100\% \\
22 & 0.680 $\pm$ 0.117 & +0.0001 & $-$0.0019 to +0.0023 & yes & yes & 100\% \\
21 & 0.678 $\pm$ 0.116 & $-$0.0021 & $-$0.0042 to $-$0.0001 & yes & yes & 100\% \\
20 & 0.678 $\pm$ 0.117 & $-$0.0027 & $-$0.0055 to +0.0002 & yes & yes & 100\% \\
19 & 0.679 $\pm$ 0.118 & $-$0.0007 & $-$0.0032 to +0.0018 & yes & yes & 100\% \\
18 & 0.677 $\pm$ 0.117 & $-$0.0031 & $-$0.0059 to $-$0.0004 & yes & yes & 100\% \\
17 & 0.675 $\pm$ 0.115 & $-$0.0057 & $-$0.0097 to $-$0.0015 & yes & yes & 100\% \\
16 & 0.672 $\pm$ 0.114 & $-$0.0081 & $-$0.0122 to $-$0.0040 & yes & yes & 100\% \\
15 & 0.667 $\pm$ 0.115 & $-$0.0136 & $-$0.0187 to $-$0.0085 & yes & no & 100\% \\
14 & 0.665 $\pm$ 0.112 & $-$0.0152 & $-$0.0214 to $-$0.0090 & no & no & 0\% \\
13 & 0.660 $\pm$ 0.110 & $-$0.0206 & $-$0.0280 to $-$0.0135 & no & no & 0\% \\
12 & 0.658 $\pm$ 0.109 & $-$0.0218 & $-$0.0298 to $-$0.0141 & no & no & 0\% \\
11 & 0.653 $\pm$ 0.110 & $-$0.0272 & $-$0.0367 to $-$0.0180 & no & no & 0\% \\
10 & 0.640 $\pm$ 0.111 & $-$0.0401 & $-$0.0507 to $-$0.0295 & no & no & 0\% \\
9 & 0.628 $\pm$ 0.111 & $-$0.0520 & $-$0.0644 to $-$0.0398 & no & no & 0\% \\
8 & 0.617 $\pm$ 0.112 & $-$0.0631 & $-$0.0744 to $-$0.0516 & no & no & 0\% \\
7 & 0.608 $\pm$ 0.111 & $-$0.0726 & $-$0.0849 to $-$0.0603 & no & no & 0\% \\
6 & 0.594 $\pm$ 0.106 & $-$0.0862 & $-$0.0998 to $-$0.0729 & no & no & 0\% \\
5 & 0.591 $\pm$ 0.104 & $-$0.0889 & $-$0.1048 to $-$0.0735 & no & no & 0\% \\
4 & 0.582 $\pm$ 0.110 & $-$0.0981 & $-$0.1141 to $-$0.0825 & no & no & 0\% \\
3 & 0.544 $\pm$ 0.102 & $-$0.1360 & $-$0.1527 to $-$0.1189 & no & no & 0\% \\
2 & 0.435 $\pm$ 0.083 & $-$0.2448 & $-$0.2682 to $-$0.2206 & no & no & 0\% \\
1 & 0.305 $\pm$ 0.062 & $-$0.3749 & $-$0.4056 to $-$0.3445 & no & no & 0\% \\
\bottomrule
\end{tabular}
\par\smallskip\parbox{0.95\linewidth}{\footnotesize \rev{Ranked minus mean of 10 random subsets (Wilcoxon, Holm over 3): $k$ = 6: +0.0352 (95\% CI +0.0230 to +0.0484), $p_{\mathrm{Holm}}$ $<$ 0.001; $k$ = 12: +0.0268 (95\% CI +0.0197 to +0.0340), $p_{\mathrm{Holm}}$ $<$ 0.001; $k$ = 18: +0.0141 (95\% CI +0.0093 to +0.0191), $p_{\mathrm{Holm}}$ $<$ 0.001.}}
\end{table}


\bibliographystyle{elsarticle-num}
\bibliography{references}

\end{document}
